%% main.tex — transmon-benchmark.6 (anvil:pub revision, iteration 6)
%%
%% Revised 2026-07-16 by pub-revise from transmon-benchmark.5 + the
%% transmon-benchmark.5.review and transmon-benchmark.5.audit critic
%% siblings. SURGICAL pass (iteration 6 of max 6; v5 scored 41/44,
%% advance:true — the framing is settled): (1) resolves the audit's one
%% critical flag by unquoting the eriksson2025automated span in §1 to a
%% paraphrase (the v5 quotation dropped the source's word "iterative");
%% (2) annotates the off-target Peak RSS cells in the CPU table with
%% their provenance (the committed [matched.off_target.*] tables record
%% wall times only); (3) aligns the §8.3 GPU-loss ratio to the
%% artifact's own 136x; (4) adds the 0.2155-vs-0.2156 anchor clarifier
%% in §9.2. No restructuring. See changelog.md for the note -> change
%% map.
%%
%% ---- v5 provenance note (historical) --------------------------------
%% Revised 2026-07-16 by pub-revise from transmon-benchmark.4 + the
%% transmon-benchmark.4.review and transmon-benchmark.4.litsearch critic
%% siblings, against the ⭐ 2026-07-16 REFRAME in
%% papers/transmon-benchmark/BRIEF.md and the operator-approved spine at
%% docs/research/transmon-paper-reframe.md (PR #587): the paper's
%% contribution is now DIFFERENTIABLE TRANSMON DESIGN (the LOM branch),
%% with the cross-validation benchmark as the correctness credential.
%%
%% Version-provenance note: transmon-benchmark.4/main.tex received an
%% out-of-band in-place edit (PR #557, 2026-07-16) that corrected the
%% 1.16M-DOF scale story in §8.2 only; the v4 review flagged six other
%% sites still carrying the retracted OOM/memory-wall claim
%% (critical flag `numerical_inconsistency_scale_story`). This revision
%% propagates the corrected flop-and-fill finding everywhere.
%%
%% Every agreement number is quoted VERBATIM from the committed
%% benchmarks/transmon_eigen/results.toml. CPU-cell numbers trace to
%% benchmarks/transmon_bench_cpu/results.toml and the committed 1.16M-DOF
%% run log benchmarks/transmon_bench_cpu/geode_runs_1p16M_2026-07-15.log.
%% GPU scaling numbers trace to benchmarks/gpu_driven_scaling/results.toml.
%% Quantum/LOM numbers trace to benchmarks/transmon_quantum/results.toml.
%% Differentiable-design numbers trace to
%% benchmarks/transmon_diffopt/results.toml (PR #588) and
%% benchmarks/transmon_diffopt/pad_results.toml (PR #590); adjoint-layer
%% validation numbers trace to the merged test suites of
%% crates/geode-core/src/{adjoint.rs, shape.rs, driven/adjoint.rs,
%% driven/shape.rs} (PRs #573/#575/#579/#581/#586). Eigensolve-roadmap
%% numbers trace to
%% benchmarks/transmon_bench_cpu/sigma4p5_deepshift_characterization.md.
%% See changelog.md in this version dir for the note -> change map.

\documentclass{anvil-paper}

%% --- URL line-breaking: the two AWS-blog references carry very long
%% hyphenated URLs; allow breaks at hyphens and add stretch so the
%% bibliography stays inside the measure (render-gate overfull check).
\expandafter\def\expandafter\UrlBreaks\expandafter{\UrlBreaks\do\-}
\Urlmuskip=0mu plus 1mu

%% --- Draft-phase helper macros -------------------------------------------
%% Figure placeholder: uses the rendered figure if pub-figures has produced
%% it, otherwise a framed placeholder naming the source script.
\newcommand{\anvilfig}[3][0.85]{%
  \IfFileExists{figures/#2}%
    {\includegraphics[width=#1\linewidth]{figures/#2}}%
    {\fbox{\parbox{#1\linewidth}{\centering\vspace{5em}%
      \texttt{figures/#2}\\[0.5em]%
      {\small placeholder --- render via \texttt{figures/src/#3}}%
      \vspace{5em}}}}}

\title{Differentiable transmon design with a tensor-native finite-element
electromagnetics solver: gradient-based optimization of charging energy,
cross-validated against Palace}
% TODO(operator): final title wording is operator-approved per BRIEF.md
% ⭐ REFRAME (2026-07-16); this is the BRIEF's suggested title class.
\author{Robb Walters\thanks{These authors contributed equally.}
  \and Crutcher Dunnavant\thanks{These authors contributed equally.}}
% TODO(operator): affiliations (author order and equal-contribution
% footnote are operator-resolved per BRIEF.md, 2026-07-14).
\date{\today}

\begin{document}
\maketitle

\begin{abstract}
Superconducting-qubit electromagnetic design is an iterative
guess-and-check loop: the workflow's production field solvers expose no
derivatives of design figures-of-merit, so state-of-practice optimizers
iterate time-consuming electromagnetic simulations and inject parameter
updates from separate analytic physics models
\citep{eriksson2025automated}. We demonstrate the missing capability ---
\emph{solver-derived} design gradients from a 3D finite-element
electromagnetics solve --- and use it for gradient-based optimization of a
transmon's charging energy $E_C$. GEODE-FEM is a Rust solver built on the
Burn tensor framework; a discrete-adjoint layer bridges the one stage its
autodiff tape cannot reach (the sparse factorization), yielding a
validated $2\times 2$ sensitivity matrix --- material ($\partial/\partial
\varepsilon$) and geometry ($\partial/\partial X$) derivatives on both the
electrostatic and the driven $H(\mathrm{curl})$ operators --- each
finite-difference-validated with mutation tests (relative errors
$10^{-9}$--$10^{-5}$). Because the capacitance $C = \phi^{\mathsf{T}}
K\phi$ is variationally stationary, its shape derivative collapses to a
Hellmann--Feynman-like explicit term; composing it with the
capacitance-to-$E_C$ chain, Newton iteration drives a capacitor geometry
to a $0.2156$\,GHz target $E_C$, confirmed by an independent fresh solve
to $1.4\times 10^{-15}$, and on the real 133k-tetrahedron device mesh
$\partial C_\Sigma/\partial\theta$ is FD-validated at $1.15\times 10^{-4}$
with a genuine two-step convergence to a within-budget target. The honest
negatives are stated, not buried: the fixed-topology pad parametrization
falls $33\times$ short of the ${\sim}90$\,fF design anchor (a
mesh-morphing problem, named as future work), and eigenmode/EPR
derivatives remain roadmap, blocked on a characterized deep-shift
eigensolve wall. The correctness credential: on the identical mesh,
GEODE-FEM reproduces all six of Palace's transmon eigenmodes to
$\leq 0.033\%$ (worst case $0.032\%$). The contribution is a capability,
not a speed race --- on raw performance the measured record shows no
corner where GEODE-FEM beats Palace at scale, and we report it honestly.
Every number traces to a committed artifact.
\end{abstract}

\section{Introduction}
\label{sec:intro}

The hard problem in superconducting-qubit electromagnetic design is the
\emph{inverse map}: given target Hamiltonian parameters --- a charging
energy $E_C$, an anharmonicity $\alpha$, coupling capacitances --- produce
a physical layout. Today that map is walked by hand: parameterize a
geometry, run a 3D electromagnetic simulation, extract the parameters,
compare, adjust, repeat. The state-of-practice automated workflow makes
the loop's structure explicit: QDesignOptimizer
\citep{eriksson2025automated} iterates time-consuming electromagnetic
simulations of HFSS-class solvers and guides the parameter updates with
\emph{separate}, user-defined analytic physics models --- precisely
because the field solver at the center of the loop produces observables
but no derivatives of those observables with respect to the design. The
gradient the optimizer actually wants, $\partial(\text{figure of
merit})/\partial(\text{geometry, materials})$ through the field solve
itself, is structurally unavailable from the production solvers ---
commercial (HFSS, COMSOL) and open-source (Palace \citep{palace}) alike:
their assembly and factorization pipelines are hand-written kernels with
no derivative path, and none exposes an adjoint of its solve.

The surrounding literature brackets this gap without closing it. At the
\emph{lumped-circuit} level, differentiation is a solved problem:
SQcircuit \citep{rajabzadeh2023analysis} analyzes arbitrary
superconducting circuits, and its gradient-based follow-up
\citep{rajabzadeh2024general} computes eigenvalue and eigenvector
gradients of the circuit Hamiltonian via PyTorch autodiff, optimizing
qubit designs in circuit-parameter space. But the map from geometry and
materials \emph{to} those lumped parameters still runs through
non-differentiable 3D electromagnetic simulation. At the \emph{field}
level, differentiable electromagnetics is owned by photonics: adjoint and
autodiff methods on FDTD and integral-equation solvers are a mature
inverse-design industry \citep{molesky2018inverse,hughes2019forward,
hammond2022high,schubert2022inverse,mahlau2026fdtdx,ponomareva2025torchgdm},
and differentiable FEM at scale has been proven outside electromagnetics
by JAX-FEM in solid mechanics \citep{xue2023jax}. What none of these
supplies is the regime superconducting-qubit (and, more broadly, RF)
design lives in: \emph{frequency-domain finite elements on body-fitted
tetrahedra with anisotropic materials} --- distributed, field-level EM
gradients from the same FEM solve that produces the design observables.
That wedge is, to our knowledge, uncontested; this paper occupies it.

GEODE-FEM (``geode-fem'') is positioned to close the gap by construction
rather than retrofit. It is a full-wave 3D $H(\mathrm{curl})$/N\'ed\'elec
(plus scalar electrostatic) finite-element solver written in Rust on the
Burn tensor framework \citep{burn}: element assembly is batched
$[n_{\mathrm{elem}},6,6]$ tensor algebra on an autodiff-capable substrate,
so the derivative of assembly with respect to materials and node
coordinates is available where a classical solver has only hand-written
kernels. One stage breaks the tape: the sparse direct factorization
(faer \citep{kazdadi2026faer}) at the heart of every solve. A
discrete-adjoint layer closes exactly that gap --- one extra
(transpose-)solve, reusing the forward factorization, restores the full
gradient (Section~\ref{sec:adjoint}). The resulting capability is
general --- shape and material sensitivities for RF, photonic, and
electrostatic device design --- and the transmon is our demonstration
vehicle because its unmet need is the clearest documented in the
literature, not because the method is qubit-specific.

A design-gradient claim from a new solver is only as credible as the
solver's forward answers, so the paper carries its correctness credential
up front: on a real transmon-plus-readout-resonator geometry meshed once
and consumed identically by both codes, geode-fem reproduces Palace ---
the MFEM-based \citep{mfem,andrej2024high} reference solver for
superconducting-qubit design --- to within $0.033\%$ across all six
eigenmodes (worst case $0.032\%$; junction LC mode $0.001\%$), with exact
analytic tripwires guarding against coincidental agreement
(Section~\ref{sec:agreement}). This parity is \emph{a credential, not the
contribution} --- matching Palace is the price of entry --- and it is,
incidentally, the first independent third-party validation of the
Palace/DeviceLayout.jl transmon stack, which has no publication of its
own.

Our contributions are:
\begin{itemize}
  \item \textbf{A validated $2\times 2$ discrete-adjoint sensitivity
    matrix for FEM electromagnetics.} Material
    ($\partial/\partial\varepsilon$) and geometry ($\partial/\partial X$,
    exact $\partial K/\partial X$ via forward-mode duals through the
    element kernels) derivatives, on both the scalar electrostatic and
    the complex driven $H(\mathrm{curl})$ N\'ed\'elec operators --- each
    one forward plus one adjoint solve sharing a single factorization,
    each validated against full-pipeline central finite differences with
    mutation tests that reject sign, conjugation, and dropped-term errors
    (Section~\ref{sec:adjoint}).
  \item \textbf{A differentiable capacitance$\to E_C$ chain with a
    variational-stationarity simplification.} Because $C =
    \phi^{\mathsf{T}} K\phi$ is stationary in the solved potential, the
    implicit term vanishes and the shape derivative is a pure
    explicit-geometry (Hellmann--Feynman-like) term; validated against
    the analytic parallel-plate law and central finite differences
    (Section~\ref{sec:chain}).
  \item \textbf{Gradient-based optimization demonstrations, honest tiers
    labeled.} An end-to-end Newton loop that hits a $0.2156$\,GHz target
    $E_C$ and is confirmed by an independent fresh solve at
    $1.4\times 10^{-15}$ (the parametrization is affine, so one-step
    convergence proves the loop, not a hard optimization); a real-device
    demonstration on the committed 133k-tet mesh with the gradient
    FD-validated at $1.15\times 10^{-4}$ and a genuine two-step
    convergence; and the honest negative --- the fixed-topology pad
    parametrization falls $33\times$ short of the ${\sim}90$\,fF design
    anchor, a mesh-morphing problem we name rather than hide
    (Section~\ref{sec:optdemo}).
  \item \textbf{The correctness credential.} Same-mesh, same-physics,
    same-junction-model cross-validation against Palace at
    $\leq 0.033\%$ across all six modes, with oracle-first methodology
    and honest-physics reporting --- the by-construction absent qubit
    mode and a disclosed, diagnosed spurious mode
    (Sections~\ref{sec:agreement} and~\ref{sec:honest-physics}) ---
    doubling as the first independent validation of the
    Palace/DeviceLayout.jl stack.
  \item \textbf{Honest performance and roadmap accounting.} A matched
    CPU cell shows a per-core, small-to-medium-scale win for the direct
    eigensolve that \emph{loses at scale on both axes} --- at
    ${\sim}1.16$M DOFs it completes given adequate memory ($565.5$\,s,
    $92.2$\,GB peak) but Palace is faster and far leaner ($423.12$\,s,
    ${\sim}33$\,GB aggregate); a flop-and-fill crossover below 1M DOFs,
    not merely a memory wall --- and the GPU cell is correctness-proven,
    performance-negative at measured sizes (Section~\ref{sec:perf}).
    Eigenmode/EPR derivatives are \emph{not yet differentiable} here;
    the blocker and the named path are stated as roadmap, with no
    overclaim (Section~\ref{sec:roadmap}).
\end{itemize}

Section~\ref{sec:related} positions the work; Section~\ref{sec:architecture}
presents the substrate and where its tape breaks;
Sections~\ref{sec:problem}--\ref{sec:honest-physics} establish the
benchmark problem and the correctness credential;
Sections~\ref{sec:lom}--\ref{sec:optdemo} are the contribution ---
the forward LOM chain, the adjoint layer, and the optimization
demonstrations; Section~\ref{sec:roadmap} is the eigenmode roadmap;
Sections~\ref{sec:perf}--\ref{sec:discussion} report performance context,
reproducibility, and discussion.

\section{Related Work}
\label{sec:related}

\paragraph{Superconducting-qubit design optimization.}
The design loop this paper attacks is documented across the qubit-design
software stack. QDesignOptimizer \citep{eriksson2025automated} is the
state of practice: physics-guided multi-parameter optimization that
wraps non-differentiable EM simulation (HFSS with pyEPR/Qiskit-Metal-class
post-processing) and steers it with analytic models, because no gradient
is available from the solver. The SQcircuit line
\citep{rajabzadeh2023analysis,rajabzadeh2024general} differentiates at
the lumped-circuit level --- including eigenvalue/eigenvector gradients
of the circuit Hamiltonian via PyTorch --- and optimizes in circuit-parameter
space; the geometry-to-parameters map remains simulation-bound. The
SQuADDS database \citep{shanto2024squadds} attacks the same loop by
lookup: pre-simulated, experimentally-validated design points. Qiskit
Metal's quasi-lumped methods \citep{minev2021circuit} and the
EPR/black-box quantization physics
\citep{koch2007,nigg2012,minev2021} frame the parameter-extraction layer
this paper's LOM chain implements. None of these obtains gradients from
the field solve itself; that is the gap Section~\ref{sec:adjoint} closes.

\paragraph{Differentiable electromagnetics --- owned by photonics.}
Adjoint/autodiff EM is a mature inverse-design discipline in
nanophotonics \citep{molesky2018inverse}, expressed through structured
grids and integral methods: autograd FDFD/FDTD in ceviche
\citep{hughes2019forward}, foundry-constrained JAX inverse design
\citep{schubert2022inverse}, the Meep adjoint pipeline
\citep{hammond2022high}, the JAX-native FDTDX framework
\citep{mahlau2026fdtdx}, GPU-accelerated autodiff integral-method
scattering in TorchGDM \citep{ponomareva2025torchgdm}, and
differentiable RCWA \citep{kim2024meent}. These codes accept rectilinear
grids or Green's-function discretizations; the superconducting-qubit
regime --- body-fitted tetrahedra, anisotropic sapphire,
$H(\mathrm{curl})$ conformity, frequency domain --- is exactly where that
literature stops. Differentiable FEM has been proven outside EM: JAX-FEM
\citep{xue2023jax} (GPU-accelerated, autodiff-native continuum mechanics)
and PyTorch-FEA \citep{liang2023pytorch} in biomechanics; concurrent work
on general differentiable Galerkin assembly (TensorGalerkin
\citep{wen2026learning}, with the torch-sla adjoint sparse-solver library
\citep{chi2026torch}) independently confirms the substrate-level thesis
while reporting nodal magnetostatics rather than full-wave
$H(\mathrm{curl})$ EM. To our knowledge, no prior work delivers
field-level design gradients from a frequency-domain 3D FEM
electromagnetics solve of the class RF/superconducting design uses.

\paragraph{FEM kernels by code generation, and tensor compilers.}
Generating FEM kernels from high-level descriptions is a two-decade
tradition --- FFC \citep{kirby2006compiler}, UFL
\citep{alnaes2014unified}, Firedrake \citep{rathgeber2017firedrake}, TSFC
\citep{homolya2018tsfc} --- whose modern GPU expression is libCEED
\citep{brown2021libceed} and MFEM's partial-assembly path
\citep{andrej2024high}. Palace itself runs on that domain-specific stack.
geode-fem's substrate bet is categorically different: outsource kernels,
scheduling, and backend portability to a \emph{general-purpose ML tensor
stack} (the Halide/TVM compiler lineage
\citep{ragankelley2013halide,chen2018tvm}, reaching scientific computing
via ML frameworks), and phrase assembly and matrix-free application as
batched tensor algebra. For this paper the decisive property of that
choice is not portability but \emph{differentiability}: an
autodiff-capable substrate makes the assembly derivative available by
construction, which is what the adjoint layer of
Section~\ref{sec:adjoint} composes with. Adjacent genera --- simulation
DSLs \citep{hu2019difftaichi,holl2020learning}, neural-ansatz
eigensolvers \citep{jiang2024fieldtnn}, hand-written CUDA FEM
\citep{huantingmeng2014gpu} --- trade away either the general-purpose
stack or the discretization itself.

\paragraph{Validation and cross-solver benchmarking.}
Code-to-code comparison on precisely specified problems is an established
genre in computational electromagnetics (the TEAM workshop tradition
since 1986\footnote{The TEAM founding documents are workshop reports
without resolvable DOIs; we describe the tradition in prose.}; nano-optics
solver comparisons \citep{hoffmann2009comparison}); in V\&V vocabulary
\citep{roache1997quantification,oberkampf2010verification} our credential
is verification-side, code-to-code evidence. On the qubit-EM axis the
closest works are SQDMetal \citep{sommers2025open} (Palace cross-checked
against COMSOL/HFSS at the $0.02$--$0.13\%$ level, independently
constructed models) and the Palace-centered workflow of
\citet{ye2025electromagnetic} (validated against cryogenic measurement at
${\sim}0.3\%$). Both wrap and trust Palace; neither is a same-mesh
comparison against an independently \emph{implemented} solver. Our
protocol eliminates meshing and model entry as confounders --- identical
MSH fixture, identical lumped-shunt junction, first-order N\'ed\'elec on
both sides --- so residuals isolate formulation and solver correctness.
Mature open EM tools normally carry a methods paper (GetDP
\citep{dular1998general}, FEniCS/DOLFIN \citep{logg2010dolfin}, Gmsh
\citep{gmsh}, scikit-rf \citep{arsenovic2022scikit}, MFEM
\citep{mfem}); Palace and DeviceLayout.jl are unpublished outliers
\citep{palace,palace-blog,devicelayout,devicelayout-blog,
devicelayout-juliacon}, which is the gap the credential narrows from the
outside. geode-fem's numerics: first-order N\'ed\'elec elements
\citep{nedelec1980}, shift-invert Lanczos \citep{arpack} over a real
symmetric pencil factorized with faer \citep{kazdadi2026faer}; Palace:
SLEPc-class Krylov \citep{slepc} with divergence-free projection and AMS
preconditioning \citep{hiptmair-xu,
tzaniovkolevpanayotsvassilevski2018parallel}; the gauge-fixing literature
behind Section~\ref{sec:honest-physics} is the tree--cotree tradition
\citep{albanese1988solution,manges1995generalized}.

\section{GEODE-FEM: A Tensor-Native Substrate, and Where the Tape Breaks}
\label{sec:architecture}

geode-fem is written in Rust on the Burn tensor framework \citep{burn},
whose cubecl JIT compiles tensor kernels for CPU, CUDA, and WebGPU
backends from one source. Every throughput-critical FEM stage is phrased
as the batched dense algebra the stack already optimizes.

\paragraph{Batched element-local assembly.}
For first-order N\'ed\'elec tetrahedra the element curl--curl and mass
contributions are $6\times 6$ matrices depending on per-element geometry
factors and material tensors (including the rotated anisotropic sapphire
permittivity of Section~\ref{sec:geometry}); the scalar electrostatic
operator is the analogous P1 kernel. geode-fem computes them for all
elements at once as batched $[n_{\mathrm{elem}},6,6]$ tensor operations
--- one kernel launch per assembly stage rather than an element loop ---
then either scatters into a global sparse matrix (the assembled f64 CPU
path) or retains the batch for matrix-free application.

\paragraph{Matrix-free application and on-device Krylov.}
The matrix-free operator apply is a three-stage tensor pipeline ---
gather global DOFs to element-local vectors, batched matmul against the
stored element operators, scatter-add back --- and the driven
frequency-sweep solve runs a conjugate orthogonal conjugate gradient
(COCG) iteration entirely on the device with a constant number of scalar
host--device synchronizations per iteration. The backend (ndarray on CPU,
CUDA on GPU) is a type parameter, not a code path.
Section~\ref{sec:perf} reports honestly how far below the
GPU-amortization threshold the current benchmark sizes sit.

\paragraph{Honest constraints.}
Two boundaries are constraints, not footnotes. The CUDA backend
(burn-cuda 0.21) is f32-only today --- the underlying cubecl runtime
currently disables f64\footnote{Upstream constraint of the burn/cubecl
stack at version 0.21 \citep{burn}; tracking-issue URL to be added at
revision. TODO(operator).} --- so every GPU result is
mixed-precision-qualified. And the eigensolve that produces the
credential comparison is factorization-bound: its inner solves are faer
sparse-LU factorizations \citep{kazdadi2026faer}, a classical serial CPU
component outside the tensor-compiler story.

\paragraph{Where the tape breaks --- and why an adjoint layer, not
autodiff-through-the-solver.}
The substrate choice has a consequence the rest of this paper builds on:
Burn's automatic differentiation reaches everything \emph{up to} the
linear solve --- material and node-coordinate dependence of the assembled
operators is differentiable by construction --- but the sparse
factorization is an external, opaque kernel that breaks the tape.
Differentiating \emph{through} a factorization by unrolling is neither
practical nor necessary: the classical discrete-adjoint identity turns
the missing derivative into one additional linear solve against the
\emph{same} factorization (Section~\ref{sec:adjoint}). The division of
labor is deliberate --- autodiff where the substrate provides it (element
kernels, via forward-mode duals), adjoint algebra where it does not (the
solve) --- and it is what makes the gradients in this paper
\emph{solver-derived} in the sense the design loop of
Section~\ref{sec:intro} lacks: the same assembly and the same
factorization that produce the observable produce its gradient.

\section{Benchmark Problem: Geometry, Mesh, and Junction Model}
\label{sec:problem}

\subsection{Geometry and mesh provenance}
\label{sec:geometry}

The device is the \texttt{SingleTransmon} example shipped with
DeviceLayout.jl v1.15.0 \citep{devicelayout}: a transmon with readout
resonator on a sapphire substrate, enclosed in a vacuum box. The sapphire
is anisotropic, with permittivity tensor
$\varepsilon = R\,\mathrm{diag}(9.3,\, 9.3,\, 11.5)\,R^{\mathsf{T}}$
rotated approximately $36.87^\circ$ in-plane. The DeviceLayout.jl
workflow emits seven named physical groups (metal, substrate, vacuum,
junction lumped-element surface, two readout ports, exterior boundary)
and drives Gmsh \citep{gmsh} to produce a tetrahedral mesh, exported as
an MSH~4.1 fixture and pinned by SHA-256 (Section~\ref{sec:repro}).

The mesh has 22{,}684 nodes and 133{,}314 tetrahedra, yielding 156{,}863
N\'ed\'elec edge unknowns, of which 133{,}108 interior degrees of freedom
remain after eliminating perfect-electric-conductor (PEC) boundary edges.
PEC is applied on all metal surfaces and the exterior box; the two
readout ports are left open (lossless), so no resistive element enters
and the eigenproblem pencil stays real symmetric.
Figure~\ref{fig:geometry} shows the geometry with physical groups
color-coded. Every measurement in this paper --- the eigenmode
credential, the electrostatic LOM chain, and the real-device gradient
demonstration --- runs on this one committed mesh (plus a uniformly
refined variant for the scale study of Section~\ref{sec:perf}).

\begin{figure}[t]
  \centering
  \anvilfig{fig1-geometry.pdf}{fig1\_geometry.py}
  \caption{Benchmark geometry: top-down projection of the committed
  surface mesh (the SHA-256-pinned MSH~4.1 fixture) for the
  DeviceLayout.jl v1.15.0 \texttt{SingleTransmon} transmon and readout
  resonator. Surface physical groups are color-coded --- the metal ground
  plane (with the coplanar-waveguide feedline, meander resonator, and
  transmon slot as etched gaps), the four-triangle junction
  \texttt{lumped\_element} surface (enlarged in the inset, actual
  triangulation shown), the two readout ports, and the exterior-boundary
  outline; the \texttt{substrate} and \texttt{vacuum} volume groups fill
  the enclosing box. Both solvers consume this identical mesh (22{,}684
  nodes; 133{,}314 tetrahedra).}
  \label{fig:geometry}
\end{figure}

\subsection{The junction model: a lumped reactive shunt}
\label{sec:junction}

The Josephson junction is modeled identically in both solvers as a lumped
reactive shunt --- a parallel $L \parallel C$ element with
$L = 14.860$\,nH and $C = 5.5$\,fF --- attached to the four-triangle
\texttt{lumped\_element} surface group $\Gamma$. Following the uniform
lumped-port construction, with port length $\ell$ and width $w$, the shunt
contributes
\begin{equation}
  K_{\mathrm{port}} = \frac{\ell}{w\,L}\, S_\Gamma
  \qquad\text{and}\qquad
  M_{\mathrm{port}} = \frac{C\,\ell}{w}\, S_\Gamma ,
  \label{eq:port}
\end{equation}
where $S_\Gamma$ is the port-surface matrix assembled on $\Gamma$. The
eigenproblem is then the real symmetric generalized pencil
\begin{equation}
  \left(K + K_{\mathrm{port}}\right) x
  \;=\; \omega^2 \left(M + M_{\mathrm{port}}\right) x .
  \label{eq:pencil}
\end{equation}
In Palace, the same physics is entered as a reactive \texttt{LumpedPort}
($L \parallel C$, no resistance) on the same surface group; the readout
ports are configured open, so Palace's pencil is likewise real. This
one-to-one correspondence is what makes the comparison same-physics as
well as same-mesh.

\section{The Correctness Credential: Same-Mesh Cross-Validation Against
Palace}
\label{sec:agreement}

A solver claiming design gradients must first prove its forward answers;
this section is that proof. Palace's own agreement with commercial tools
is established \citep{sommers2025open}, so reproducing Palace is the bar
a new implementation must clear before anything it computes \emph{about}
its solutions (such as their derivatives) is worth discussing.

\paragraph{Protocol.}
The benchmark is oracle-first. The mesh is generated once, committed,
pinned by SHA-256, and consumed by both solvers at matched (first) order
with the identical junction model; residual disagreement is therefore
attributable to formulation or implementation, not meshing. Two exact
analytic tripwires guard against coincidental or
self-consistent-but-wrong agreement: the isolated junction LC resonance
$f_{LC} = 1/(2\pi\sqrt{LC}) = 17.60$\,GHz anchors the junction mode, and
a Josephson $\sqrt{L}$-scaling test requires that doubling $L$ move the
junction mode by exactly $1/\sqrt{2}$ --- measured: $17.4901 \to
12.37$\,GHz, committed ratio $0.7071$ (computed from the unrounded
values; Figure~\ref{fig:lscaling}). The suite also contains inverse tests
that \emph{must fail} (perturbed configurations required to disagree with
the oracle --- a pipeline that still passes under them is not measuring
anything), and the Palace oracle is deterministic: a rerun from the
committed configuration reproduced its committed \texttt{eig.csv}
bit-for-bit.

\paragraph{Mode identification.}
Modes are identified, not eyeballed. Within geode-fem, the
stiffness-participation ratio
\begin{equation}
  p \;=\; \frac{x^{\mathsf{T}} K_{\mathrm{port}}\, x}
               {x^{\mathsf{T}} \left(K + K_{\mathrm{port}}\right) x}
  \label{eq:participation}
\end{equation}
cleanly separates the junction LC mode ($p = 1.000$) from resonator and
box modes ($p = 0.000$ at the artifact's three-decimal precision). Palace
reports a field-based per-port energy participation, a differently
normalized quantity that does \emph{not} rank the modes the same way on
the committed run (all six of its port-EPR values are $\leq
5\times10^{-4}$, with the junction mode's the smallest in magnitude); the
two diagnostics are complementary, not a shared discriminant, and we do
not compare them numerically. Cross-solver mode identification therefore
rests on frequency agreement against the analytic anchor.

\begin{table}[t]
  \centering
  \caption{Eigenmode agreement on the identical mesh (22{,}684 nodes /
  133{,}314 tets / 156{,}863 N\'ed\'elec DOFs, 133{,}108 interior after
  PEC), first order in both solvers, junction as a lumped reactive shunt
  ($L = 14.860$\,nH, $C = 5.5$\,fF). geode-fem at commit
  \texttt{3174015}; Palace at commit \texttt{fba6a5b}. All values are
  quoted verbatim from the committed
  \texttt{benchmarks/transmon\_eigen/results.toml}: geode-fem frequencies
  at the committed three decimals, Palace at full precision, and the
  committed per-mode \texttt{rel\_err\_pct}, computed against the rounded
  geode-fem values (precision disclosure in Section~\ref{sec:repro}).}
  \label{tab:agreement}
  \small
  \begin{tabular}{@{}llccc@{}}
    \toprule
    Mode & & geode-fem (GHz) & Palace (GHz) & $\Delta$ (\%) \\
    \midrule
    \texttt{resonator} & readout resonator & 5.153 & 5.151335830348 & 0.032 \\
    \texttt{mode\_2} & box/cavity & 15.465 & 15.46052107794 & 0.029 \\
    \texttt{junction\_lc} & junction LC & 17.490 & 17.49010903536 & 0.001 \\
    \texttt{mode\_4} & box/cavity & 18.693 & 18.69165792915 & 0.007 \\
    \texttt{mode\_5} & box/cavity & 20.703 & 20.69755679425 & 0.026 \\
    \texttt{mode\_6} & box/cavity & 26.088 & 26.08089940472 & 0.027 \\
    \bottomrule
  \end{tabular}
\end{table}

\begin{figure}[t]
  \centering
  \anvilfig{fig2-agreement.pdf}{fig2\_agreement.py}
  \caption{Mode-frequency agreement: geode-fem versus Palace on the
  identical mesh, with per-mode relative deviations annotated. Data from
  the committed benchmark artifacts
  (\texttt{benchmarks/transmon\_eigen/results.toml} and Palace's
  committed \texttt{eig.csv}).}
  \label{fig:agreement}
\end{figure}

\begin{figure}[t]
  \centering
  \anvilfig{fig3-lscaling.pdf}{fig3\_l\_scaling.py}
  \caption{Josephson inductance-scaling tripwire: junction-mode frequency
  versus junction inductance $L$ on log-log axes. Doubling $L$ moves the
  junction mode from $17.4901$ to $12.37$\,GHz --- committed ratio
  $0.7071 = 1/\sqrt{2}$ --- exactly as $f \propto 1/\sqrt{L}$ requires.
  Measured values from the committed
  \texttt{benchmarks/transmon\_eigen/results.toml} tripwire block.}
  \label{fig:lscaling}
\end{figure}

\paragraph{Result.}
Table~\ref{tab:agreement}: all six modes agree to $\leq 0.033\%$ (worst
case $0.032\%$), with the junction LC mode --- the mode the junction
model exists to produce --- at $0.001\%$. The junction mode sits $0.6\%$
below the isolated-circuit anchor $f_{LC} = 17.60$\,GHz in \emph{both}
solvers (the small shift reflects its embedding in the field problem),
and the agreement is between genuinely independent code paths --- a
batched-tensor assembly feeding a direct-factorization Lanczos in Rust,
and a projected, preconditioned Krylov method in C++ --- so a common
implementation error is an unlikely explanation for the residual, and
the tripwires exclude self-consistent-but-wrong junction physics.

\section{Honest Physics: The Absent Qubit Mode and a Disclosed Spurious
Mode}
\label{sec:honest-physics}

\subsection{No \texorpdfstring{${\sim}4$\,GHz}{\textasciitilde 4 GHz} qubit mode exists, by construction}
\label{sec:absent-mode}

The DeviceLayout.jl blog workflow \citep{devicelayout-blog} reports a
spectrum running from $4.14$ to $5.591$\,GHz --- including a
${\sim}4$\,GHz qubit mode --- while Table~\ref{tab:agreement} contains
no mode below $5.15$\,GHz. This is a property of the model, stated plainly: the
physical transmon frequency arises from the Josephson inductance
resonating against the \emph{pad (shunt) capacitance} of roughly
$80$--$100$\,fF \citep{koch2007}, and the lumped-shunt model of
Section~\ref{sec:junction} gives the junction only its own $5.5$\,fF, so
its resonance lands at $17.5$\,GHz. Both solvers agree on the absence:
Palace's projected eigensolve, instructed to hunt at $\sigma =
4.5$\,GHz, finds nothing below the $5.15$\,GHz resonator. The blog
spectrum is not reproduced by this model, and no agreement with it is
claimed. The route to qubit quantities from this geometry is the
electrostatic LOM chain of Section~\ref{sec:lom} --- which is exactly the
chain this paper makes differentiable.

\subsection{One spurious mode: disclosed, diagnosed, characterized}
\label{sec:spurious-mode}

Honest cross-validation requires reporting where the solvers
\emph{disagree}, and there is exactly one such place. geode-fem's
headline Lanczos operates on the raw pencil of Eq.~\ref{eq:pencil}
without a gauge projection, and on this problem it leaks one spurious
mode at $3.4528$\,GHz, junction-localized ($p = 0.9942$), with no Palace
counterpart. The artifact was run to ground through a measured
three-step gauge/projection arc, every step committed
(Section~\ref{sec:repro}): (1) classical tree--cotree elimination
\citep{albanese1988solution,manges1995generalized} removes it but is not
spectrum-preserving for the generalized eigenproblem --- the resonator
shifts $1.64\%$ ($5.1528 \to 5.2372$\,GHz), outside our bar, pinned as a
regression test; (2) the spectrum-preserving bulk divergence-free
projector (structurally Palace's confinement) preserves the cavity
spectrum ($0.029\%$) and collapses the 13{,}747-mode gradient cluster,
but \emph{deflates the physical junction mode} --- a lumped inductor is a
quasi-static, curl-free flux path (measured divergence ratio $50.2$;
projected norm $1.06\times10^{-4}$) --- while the spurious mode, being
genuinely solenoidal (divergence ratio ${\approx}6\times10^{-15}$),
\emph{survives}; (3) a port-aware rank-1 re-admission of the junction
flux direction yields a projector that retains all six physical modes at
$\leq 0.029\%$ worst case versus Palace. What the artifact \emph{is}:
under the L-doubling harness it moves $3.4528 \to 2.4449$\,GHz (ratio
$0.7081$, the junction $1/\sqrt{L}$ law) with $99.4\%$ of its stiffness
energy in the $K_{\mathrm{port}}$ surface term --- a second LC-like
resonance of the port-surface operator $S_\Gamma$, which Palace's
\texttt{LumpedPort} elimination represents differently and therefore
never sees. The headline comparison deliberately remains the committed
ungauged benchmark (the port-aware solver shipped after the headline
measurement); the filter now excludes a \emph{characterized} artifact
rather than an unexplained one. Figure~\ref{fig:participation} contrasts
the spectra. The arc measures precisely what Palace's divergence-free
projection buys in a production eigensolver, and what a lumped-element
formulation demands beyond it.

\begin{figure}[t]
  \centering
  \anvilfig{fig6-participation.pdf}{fig6\_participation.py}
  \caption{geode-fem junction-participation spectrum on the identical
  mesh, with Palace's mode locations (from the committed
  \texttt{eig.csv}) overlaid. The physical junction LC mode appears at
  $17.49$\,GHz in both solvers' spectra (geode-fem stiffness
  participation $p = 1.000$); geode-fem's un-projected headline Lanczos
  additionally leaks one spurious junction-localized mode at
  $3.4528$\,GHz ($p = 0.9942$) with no Palace counterpart ---
  characterized in Section~\ref{sec:spurious-mode} as a port-surface
  artifact obeying the junction $1/\sqrt{L}$ law. Palace's field-based
  port-EPR is a differently normalized diagnostic and is not plotted on
  geode-fem's participation axis (Section~\ref{sec:agreement}).}
  \label{fig:participation}
\end{figure}

\section{The Forward LOM Chain: From Fields to \texorpdfstring{$E_C$}{E\_C}}
\label{sec:lom}

The lumped-oscillator-model (LOM) branch computes the transmon's
Hamiltonian parameters from the \emph{electrostatic} problem on the same
committed mesh --- and it is the chain the remainder of this paper
differentiates. The mesh's metal is split into three node-disjoint
conductors (ground+resonator, feedline, island); a tensor-$\varepsilon$
electrostatic solve yields the Maxwell capacitance matrix by the
energy-reaction form $C_{ij} = \phi_i^{\mathsf{T}} K \phi_j$; the island
self-capacitance is reduced for the floating feedline
($C_\Sigma = C_{ii} - C_{if}^2/C_{ff}$); and the charging energy follows
as $E_C/h = e^2/(2 C_\Sigma h)$, with the transmon spectrum evaluated
charge-basis-exactly after \citet{koch2007}.

Measured on the committed fixture
(\path{benchmarks/transmon_quantum/results.toml}): $C_\Sigma =
136.7$\,fF, hence $E_C = 0.142$\,GHz, $E_J/E_C = 77.6$ (with $E_J$ from
the junction inductance), $\omega_{01} = 3.38$\,GHz, and anharmonicity
$\alpha = -0.158$\,GHz --- a plausible transmon computed end-to-end from
committed geometry, with the $\alpha \approx -E_C$ sanity holding.

\paragraph{The honest negative that travels with this chain.}
The extracted $C_\Sigma = 136.7$\,fF \emph{exceeds} the ${\sim}90$\,fF
design-anchor value back-solved from the blog workflow's expected
$E_C = 0.2156$\,GHz. The committed benchmark diagnoses this as a
model-definition difference, not a truncation artifact (grounding the
exterior wall moves $C_\Sigma$ by a relative $6\times10^{-5}$; the
scalar-versus-tensor sapphire treatment shifts it $0.75\%$). We report
the extracted numbers and the discrepancy; we do not retrofit to the
anchor. Section~\ref{sec:optdemo} returns to this gap as an optimization
target --- and reports honestly how far a fixed-topology parametrization
gets toward closing it.

\section{Differentiable Sensitivities: The Discrete-Adjoint Layer}
\label{sec:adjoint}

\subsection{The \texorpdfstring{$2\times 2$}{2x2} sensitivity matrix}
\label{sec:adjoint-matrix}

For a solve $A(\varepsilon, X)\,x = b(X)$ and a scalar observable
$g(x)$, the discrete adjoint $A^{\mathsf{T}}\lambda = \partial g/\partial
x$ turns the parameter gradient into explicit contractions:
\begin{equation}
  \frac{dg}{d\varepsilon_k}
    = -\lambda^{\mathsf{T}}\frac{\partial A}{\partial \varepsilon_k}x,
  \qquad
  \frac{dg}{dX}
    = \lambda^{\mathsf{T}}\frac{\partial b}{\partial X}
      - \lambda^{\mathsf{T}}\frac{\partial A}{\partial X}x
  \label{eq:adjoint}
\end{equation}
(real case; the complex driven case takes $2\,\mathrm{Re}[\cdot]$ of the
same contractions with the \emph{transpose}, not conjugate-transpose,
solve --- the pencil is complex-symmetric, so the adjoint solve reuses
the forward LU factorization). Everything on the right-hand side is
assembly-level and therefore differentiable on the tensor substrate: the
element Jacobians $\partial A/\partial X$ and $\partial b/\partial X$ are
\emph{exact}, computed by forward-mode dual numbers through the same
closed-form element kernels the solver assembles --- no finite-difference
truncation enters the gradient. Each gradient costs one forward plus one
adjoint solve sharing a single factorization (asserted programmatically:
\texttt{n\_factorizations} $= 1$).

Four adjoint modules are merged and validated, spanning material and
geometry sensitivities on both operators
(Table~\ref{tab:adjoint-matrix}). Each carries a committed test that the
analytic gradient match a central finite difference \emph{of the full
independent pipeline} (perturb the parameter, re-assemble, re-solve
through the public solver path, recompute $g$) --- and, load-bearingly,
mutation tests proving the tolerances bite: the classic complex-adjoint
conjugation error ($A^{\mathsf{H}}$ for $A^{\mathsf{T}}$) is rejected by
the FD check, as is a dropped term. The driven-shape case surfaced a
physics subtlety worth recording: the current-source right-hand side
$b = \mathrm{i}\omega\mu_0\!\int\! N\cdot J$ is geometry-dependent, so
the shape gradient carries a $\partial b/\partial X$ term absent from the
material case --- omitting it fails the finite-difference cross-check
outright.

\begin{table}[t]
  \centering
  \caption{The merged $2\times 2$ discrete-adjoint sensitivity matrix.
  Each cell: one forward + one adjoint solve sharing a single
  factorization; element Jacobians exact (forward-mode duals through the
  production kernels); validated against central finite differences of
  the full pipeline, with mutation tests. ``FD rel-err'' is the observed
  agreement of the analytic gradient with the FD reference in the merged
  test suite (test acceptance bounds are looser; see the module sources).
  Modules in \texttt{crates/geode-core/src/}.}
  \label{tab:adjoint-matrix}
  \footnotesize
  \setlength{\tabcolsep}{4pt}
  \begin{tabular}{@{}llll@{}}
    \toprule
    Operator & Parameter & Module (issue / PR) & FD rel-err \\
    \midrule
    scalar electrostatic (SPD) & material $\partial/\partial\varepsilon$ &
      \texttt{adjoint} (\#570 / \#573) & ${\sim}3\times10^{-8}$ \\
    scalar electrostatic (P1) & geometry $\partial/\partial X$ &
      \texttt{shape} (\#571 / \#575) & ${\sim}1\times10^{-9}$ \\
    driven $H(\mathrm{curl})$ & material $\partial/\partial\varepsilon$ &
      \texttt{driven::adjoint} (\#576 / \#579) & $2.3\times10^{-5}$ (worst region) \\
    driven $H(\mathrm{curl})$ & geometry $\partial/\partial X$ &
      \texttt{driven::shape} (\#577 / \#581) & $2.2\times10^{-9}$, $1.2\times10^{-8}$ (two maps) \\
    \bottomrule
  \end{tabular}
\end{table}

The matrix is deliberately $2\times 2$: material and geometry are the two
design-parameter classes, and the scalar and driven operators are the
electrostatic (LOM) and RF (S-parameter/EPR-workhorse) branches of the
solver. The LOM optimization of Section~\ref{sec:optdemo} exercises the
scalar column end-to-end; the driven column is merged and validated at
the operator level and composes with the driven solve that geode-fem's
GPU path accelerates --- RF shape optimization is the natural follow-on,
not exercised here.

\subsection{The \texorpdfstring{capacitance$\to E_C$}{capacitance-to-E\_C} chain, and a vanishing adjoint}
\label{sec:chain}

Composing the shape adjoint with the LOM chain of Section~\ref{sec:lom}
gives $\partial(C, E_C, \alpha)/\partial(\text{geometry})$ (merged as
\texttt{shape::capacitance\_shape\_gradient}, \#583 / PR \#586). The
composition carries an elegant simplification. The capacitance observable
is the stored field energy, $C = 2W = \phi^{\mathsf{T}}K(X)\phi$, and $W$
is \emph{variationally stationary} in the solved potential $\phi$: at the
forward solution the free rows of the cotangent $K\phi$ vanish
(equilibrium), so the adjoint solution $\lambda \approx 0$ and the entire
shape derivative collapses to the explicit-geometry term
\begin{equation}
  \frac{dC}{dX}
  \;=\; \phi^{\mathsf{T}}\,\frac{\partial K}{\partial X}\,\phi
  \qquad\text{(at the solution)},
  \label{eq:hf}
\end{equation}
a Hellmann--Feynman-like result: no adjoint back-solve is needed for the
capacitance itself. The implementation keeps the general adjoint form
(the mutation test measures the implicit part at $<10^{-8}$ of the
total, confirming the collapse numerically rather than assuming it), and
chains through $\partial E_C/\partial C_\Sigma = -e^2/(2C_\Sigma^2 h)$
exactly. Validation is double: against the analytic parallel-plate law
($\partial C/\partial d = -\varepsilon_0\varepsilon_r$ per unit area
under a gap scaling; $\partial C/\partial A = +2\varepsilon_0
\varepsilon_r$ under an area scaling --- opposite signs, both matched)
and against central finite differences of the full pipeline, with
observed relative errors at the $10^{-9}$ level.

\section{Gradient-Based Optimization: Demonstrations and an Honest
Negative}
\label{sec:optdemo}

Two demonstrations, in deliberately labeled honesty tiers: a closed-loop
proof on an analytic fixture, then the real device mesh. Both use damped
Newton on $E_C(\theta) - E_C^{\mathrm{target}}$ (equivalently
$C(\theta)$) with the analytic $dE_C/d\theta$ from
Section~\ref{sec:chain}; each Newton step costs one forward plus one
adjoint solve on a single factorization. The framing claim is
capability and step-count, not wall clock: a derivative-free loop
(the parameter sweeps of current practice) spends $N_{\mathrm{params}}$
extra forward solves per step to estimate what the adjoint delivers
exactly for one. No timing comparison against HFSS/Palace-based
optimizers is measured or claimed.

\subsection{Loop proof: Newton to a target \texorpdfstring{$E_C$}{E\_C} on a parallel-plate
fixture}
\label{sec:optdemo-plate}

On a parallel-plate capacitor fixture ($\varepsilon_r = 3$; 343 nodes,
1{,}296 tets; $\theta$ scales the gap, so $d = L(1+\theta)$), Newton
iteration drives the design from $C_0 = 120.0$\,fF ($E_C =
0.1614$\,GHz) to the design-anchor target $E_C = 0.2156$\,GHz. Full
Newton converges in \emph{one} step --- and we state the reason rather
than dress it up: $C \propto 1/(1+\theta)$ makes $E_C(\theta)$
\emph{affine} in $\theta$, so one-step convergence is expected, and this
tier proves the loop end-to-end (gradient, chain rule, update,
convergence test), not a hard optimization. The converged $\theta =
0.335657909$ matches the closed-form optimum to $3\times10^{-15}$. The
honesty check that matters: a \emph{fresh, independent} forward solve at
the converged geometry --- assemble, solve, extract, no adjoint machinery
--- reproduces the target $E_C$ to a relative error of
$1.4\times10^{-15}$. A damped run ($\alpha = 0.5$) converges in 13 steps
and supplies the descent illustration
(Figure~\ref{fig:diffopt}a). At the hit target, the Koch-exact spectrum
gives $E_J/E_C = 51.0$, $\omega_{01} = 4.128$\,GHz, $\alpha =
-0.247$\,GHz --- the ${\sim}4$\,GHz qubit regime the design anchor
implies. All trajectory values are committed in
\path{benchmarks/transmon_diffopt/results.toml}.

\begin{figure}[t]
  \centering
  \anvilfig[0.98]{fig5-diffopt.pdf}{fig5\_diffopt.py}
  \caption{Gradient-based optimization to a target $E_C$ via the
  discrete-adjoint shape gradient. (a) Parallel-plate loop proof:
  $|E_C - E_C^{\mathrm{target}}|$ versus Newton iteration; full Newton
  converges in one step (the parametrization is affine --- expected, and
  disclosed), the damped $\alpha = 0.5$ run supplies the 13-step descent
  curve; the converged design is confirmed by an independent fresh
  forward solve at relative error $1.4\times10^{-15}$
  (\texttt{benchmarks/transmon\_diffopt/results.toml}). (b) The real
  133k-tet DeviceLayout mesh under the fixed-topology island-scale
  parameter $\theta$: a within-budget demonstration target
  ($C_\Sigma = 137.0$\,fF) converges in two genuine (non-affine) Newton
  steps, fresh-solve confirmed at $5.6\times10^{-6}$; the bounded run
  toward the $89.9$\,fF design anchor stalls honestly at the mesh
  validity budget $\theta_{\mathrm{safe}} = -0.0073$ (first tetrahedron
  inversion at $\theta = -0.0097$), $33\times$ short of the
  $\theta \approx -0.241$ the anchor requires
  (\texttt{benchmarks/transmon\_diffopt/pad\_results.toml}).}
  \label{fig:diffopt}
\end{figure}

\subsection{The real device: a validated gradient, a genuine
convergence, and a \texorpdfstring{$33\times$}{33x} shortfall}
\label{sec:optdemo-pad}

The second tier runs on the committed 133k-tet SingleTransmon mesh
itself, with a fixed-topology shape parameter $\theta$ that scales the
island conductor in-plane about its centroid (101 island nodes move;
the mesh connectivity is fixed). Three results, in descending order of
comfort
(\path{benchmarks/transmon_diffopt/pad_results.toml}):

\paragraph{The gradient is real on the real mesh.} At $\theta = 0$ the
adjoint gradient $\partial C_\Sigma/\partial\theta = 198.198$\,fF is
validated against a central finite difference of the \emph{full
independent pipeline} (move the island nodes, re-assemble, re-solve,
re-extract): relative error $1.15\times10^{-4}$ at $h = 10^{-4}$, with a
clean $O(h^2)$ truncation sweep ($1.2\times10^{-2} \to 1.0\times10^{-3}
\to 1.15\times10^{-4} \to 1.0\times10^{-5}$ as $h$ falls
$10^{-3} \to 3\times10^{-5}$). This is the paper's load-bearing
real-geometry validation: the same machinery that passed at
$10^{-9}$-level on analytic fixtures holds up on production geometry.

\paragraph{The loop genuinely converges on the real mesh.} A
demonstration target within the mesh's validity budget ($C_\Sigma =
137.0$\,fF) converges in two Newton steps --- genuinely non-affine now
(the gradient changes by ${\sim}21\%$ along the path) --- and the
converged design
is confirmed by a fresh, independent multi-conductor extraction at
relative error $5.6\times10^{-6}$ (Figure~\ref{fig:diffopt}b). The
entire real-device run (validation sweep plus both optimization runs)
takes $10.7$\,s.

\paragraph{The honest negative: the anchor is out of this
parametrization's reach.} The natural target is the ${\sim}90$\,fF
anchor of Section~\ref{sec:lom} ($89.9$\,fF; $E_C = 0.2155$\,GHz ---
the last digit differs from the $0.2156$\,GHz quoted elsewhere because
this run targets the rounded $89.9$\,fF rather than the back-solved
$89.843$\,fF; both values are artifact-faithful). The
$\theta = 0$ gradient puts it at $\theta \approx -0.241$ --- but the
fixed-topology map inverts its first tetrahedron at $\theta = -0.0097$,
and the run's conservative budget (minimum tet-volume ratio $0.25$) caps
$\theta_{\mathrm{safe}} = -0.0073$, where the bounded Newton run stalls
at $C_\Sigma = 136.5$\,fF: a $33\times$ shortfall. The diagnosis is
geometric, not a gradient failure: the island's junction-attachment
nodes sit ${\sim}225\,\mu$m from the island centroid, so the centroid
scale drags them against ${\sim}0.7\,\mu$m junction-region tetrahedra.
Closing the anchor gap is therefore a \emph{mesh-morphing and
parametrization} problem --- a smoothed velocity field or per-step
remeshing, multi-parameter pad/gap/lead controls that remove the
junction lever arm, and a tensor-$\varepsilon$ extension of the shape
chain (the committed run uses the scalar trace-averaged sapphire; the
scalar-versus-tensor difference is quantified at $0.75\%$ in the
artifact) --- all named follow-ons, none of which the present paper
claims. We regard this shortfall as a result, not an embarrassment: it
is precisely the boundary information a design-gradient capability must
publish before anyone builds an optimizer on it.

\section{Roadmap: Differentiating the Eigenmode Branch}
\label{sec:roadmap}

The LOM branch above differentiates an \emph{electrostatic} solve. The
transmon's dynamical figures-of-merit --- resonator frequency,
participation-ratio and EPR self/cross-Kerr \citep{minev2021,nigg2012}
--- live in the \emph{eigenmode} problem, and that branch is \textbf{not
differentiable yet} in geode-fem. The paper commits to the explicit
split:

\begin{table}[h]
  \centering
  \small
  \begin{tabular}{@{}lll@{}}
    \toprule
    Branch & Observables & Differentiable? \\
    \midrule
    LOM (electrostatic) & $E_C$, $\alpha \approx -E_C$, coupling $C$ &
      \textbf{yes} (Sections~\ref{sec:adjoint}--\ref{sec:optdemo}) \\
    eigenmode / EPR & resonator $f$, participation, Kerr &
      \textbf{not yet} (this section) \\
    \bottomrule
  \end{tabular}
\end{table}

Differentiating an eigenpair needs the adjoint-eigenpair
(Hellmann--Feynman / Nelson) formulas
\citep{nelson1976simplified}, which require solving the interior
eigenproblem robustly at the physical shift --- and geode-fem's
matrix-free interior eigensolve is blocked at exactly that shift. The
blocker is measured and committed
(\path{benchmarks/transmon_bench_cpu/sigma4p5_deepshift_characterization.md}):
at $\sigma = 4.5$\,GHz on the 133k-DOF fixture, a \emph{single} inner
AMS-preconditioned MINRES solve (outer Lanczos step 0) never converges
--- it plateaus on a nearly flat residual tail below ${\sim}10^{-5}$ and
was still running at 14{,}300 inner iterations ($\lVert r\rVert \approx
2.6\times10^{-6}$) when killed. The plateau is
\emph{coarse-solve-invariant} --- an exact coarse solve stalls the same
way --- so the limiter is the SPD-proxy preconditioner $K + |\sigma|M$
itself, not the coarse solve. Nor is there a single inner tolerance that
is both cheap and correct: at an inner tolerance of $10^{-2}$ the entire
eigensolve completes in ${\sim}56$\,s but returns the non-physical
$\lambda \approx 0$ gradient near-kernel cluster ($0.42$--$0.64$\,GHz),
not the physical band. The named path --- not built, cited as method,
not result --- is a preconditioned Helmholtz-projected Jacobi--Davidson
iteration, applying the gradient-kernel projection inside the correction
equation, for which \citet{liang2026adaptive} provide the current
adaptive-multilevel treatment of exactly this singular Maxwell
eigenproblem; combined with the Nelson/Hellmann--Feynman eigenpair
derivatives, that is the route to differentiable eigenmode/EPR
quantities. Until it lands, no claim in this paper extends to
derivatives of the eigenmode spectrum.

\section{Performance Context}
\label{sec:perf}

Performance is context here, not contribution: on the measured record
there is no corner where geode-fem beats Palace at scale, and the
capability claim of this paper does not depend on one. We report the
matched CPU cell (with the corrected large-scale finding) and the GPU
cell (correctness proven; scaling an honest negative) for completeness
and because the honest accounting bounds what the differentiable solver
costs to run.

\subsection{CPU cell: a matched head-to-head, and a flop-and-fill
crossover}
\label{sec:perf-cpu}

All CPU measurements ran on a single EC2 \texttt{m6i.4xlarge} (8
physical cores / 16 vCPUs, ${\sim}61$\,GB usable) in
\texttt{us-west-2}.\footnote{On-demand cost approximately \$0.77/hour at
the time of the runs; the CPU cell is reproducible for a few dollars.}
The cell is matched: same box, same session, identical mesh and
eigenproblem, both solvers at 1 and 8 cores/ranks, three workloads
(physical target: 6 modes at 4.5\,GHz; off-target: 12 modes at 20\,GHz;
uniformly refined ${\sim}1.16$M-DOF mesh), full pipeline, mean of $n=3$
runs (\path{benchmarks/transmon_bench_cpu/results.toml}).

\begin{table}[t]
  \centering
  \caption{Matched CPU-cell wall-clock and memory, full pipeline,
  \texttt{m6i.4xlarge}, same box and session (mean of $n=3$).
  geode-fem \texttt{@3174015}: serial/threaded faer sparse-LU
  shift-invert Lanczos (direct); Palace \texttt{@fba6a5b}: distributed
  MPI Krylov + AMS (iterative). Palace memory is per-rank maximum RSS,
  not an aggregate. At the physical target geode-fem on \emph{one} core
  beats Palace on \emph{eight} ranks in absolute wall clock; at
  ${\sim}1.16$M DOFs the direct solve \emph{completes given adequate
  memory but loses to Palace on both wall clock and memory} --- the
  crossover is a flop-and-fill deficit, not merely a memory wall. Values
  from the committed
  \texttt{benchmarks/transmon\_bench\_cpu/results.toml}
  (\texttt{[matched.*]} tables) and, for the 128\,GB-box completion row,
  the committed run log
  \texttt{benchmarks/transmon\_bench\_cpu/geode\_runs\_1p16M\_2026-07-15.log}.}
  \label{tab:cpu}
  \small
  \begin{tabular}{@{}llcc@{}}
    \toprule
    Workload & Solver (cores) & Wall (s) & Peak RSS \\
    \midrule
    \multicolumn{4}{@{}l}{\emph{Physical target: 6 modes @ 4.5\,GHz,
      133{,}108 interior DOFs}} \\
    & geode-fem (1) & 28.7 & 3.1\,GB \\
    & geode-fem (8) & 29.0 & 3.1\,GB \\
    & Palace (1)    & 130.9 & 0.5\,GB/rank \\
    & Palace (8)    & 44.5 & 0.5\,GB/rank \\
    \addlinespace
    \multicolumn{4}{@{}l}{\emph{Off-target: 12 modes @ 20\,GHz,
      133{,}108 interior DOFs}} \\
    & geode-fem (1) & 36.8 & 3.1\,GB\textsuperscript{\ddag} \\
    & geode-fem (8) & 26.6 & 3.1\,GB\textsuperscript{\ddag} \\
    & Palace (1)    & 248.0 & 0.5\,GB/rank\textsuperscript{\ddag} \\
    & Palace (8)    & 64.7 & 0.5\,GB/rank\textsuperscript{\ddag} \\
    \addlinespace
    \multicolumn{4}{@{}l}{\emph{Large scale: 6 modes @ 4.5\,GHz,
      1{,}157{,}564 interior DOFs}} \\
    & geode-fem (1), ${\sim}61$\,GB box\textsuperscript{\dag} &
      killed at ceiling & 63.9\,GB (truncated) \\
    & geode-fem (1), 128\,GB box\textsuperscript{\dag} & 565.5 & 92.2\,GB \\
    & Palace (8)    & 423.12 & 4.1\,GB/rank (${\sim}33$\,GB aggr.) \\
    \bottomrule
  \end{tabular}

  \medskip
  \begin{minipage}{0.92\linewidth}
  \footnotesize
  \textsuperscript{\dag}The geode-fem large-scale rows depart from the
  same-box protocol of the matched cell and are marked accordingly: the
  ${\sim}61$\,GB-box run was SIGKILL-ed mid-factorization at the memory
  ceiling (its 63.9\,GB ``peak'' is a truncation of the true footprint,
  not a measurement of it), and the completion row was re-run on a
  128\,GB \texttt{r6i}-class box (committed log above), where the true
  peak is 92.2\,GB. Palace's row is from the original matched session.

  \textsuperscript{\ddag}The committed \texttt{[matched.off\_target.*]}
  tables record wall times only; the off-target Peak RSS cells repeat
  the physical-target measurements of the same session and the same
  133{,}108-DOF pencil, and are not independently recorded in the
  artifact.
  \end{minipage}
\end{table}

\begin{figure}[t]
  \centering
  \anvilfig{fig4-cpu-wallclock.pdf}{fig4\_cpu\_wallclock.py}
  \caption{Matched CPU-cell wall-clock comparison at the physical target
  (6 modes, 4.5\,GHz, 133k DOFs): geode-fem at 1 and 8 threads versus
  Palace at 1 and 8 MPI ranks. The inset reports per-core cost as
  core-seconds --- geode-fem on one core consumes $28.7$ core-s against
  Palace's $356.0$ at 8 ranks (${\sim}12\times$). Full-pipeline timings
  from Table~\ref{tab:cpu}.}
  \label{fig:cpu}
\end{figure}

Three readings, none of them a parallelization claim. \emph{Per-core
efficiency:} at the physical target, geode-fem's serial direct
factorization on one core ($28.7$\,s) beats Palace on eight ranks
($44.5$\,s) in absolute wall clock at ${\sim}12\times$ fewer
core-seconds ($28.7$ versus $356.0$); geode-fem's own 8-thread run gives
essentially no speedup ($29.0$\,s; tracked as issue \#518), so the
advantage is direct-versus-iterative throughput, not parallelism.
\emph{Target robustness:} the direct shift-invert is target-insensitive
while Palace's iterative solve degrades off-target, so the gap
\emph{widens} at 20\,GHz --- $6.7\times$ serial ($248.0/36.8$),
$2.4\times$ at 8-wide ($64.7/26.6$). \emph{The corrected large-scale
finding:} the direct factorization trades memory and super-linear
fill/flops for that speed and robustness, and the trade stops paying
below 1M DOFs. On the ${\sim}1.16$M-DOF mesh, given adequate memory
(a 128\,GB box), the direct solve \emph{does complete} --- setup
$35.78$\,s, solve $529.75$\,s, total $565.5$\,s at $92.2$\,GB peak RSS
--- but it \emph{loses to Palace on both axes}: Palace finishes the same
mesh in $423.12$\,s at $4.1$\,GB/rank (${\sim}33$\,GB aggregate). An
earlier ``never completes'' reading (the OOM kill at $63.9$\,GB on the
${\sim}61$\,GB box) was a small-box truncation artifact and is retracted
in favor of this measurement: the crossover is a \emph{flop-and-fill
deficit, not merely a memory wall} --- even with abundant memory, the
direct path is a small-to-medium-scale win that loses at scale. The
architecture trade-off, stated symmetrically: geode-fem wins
small-to-medium on speed, per-core efficiency, and target robustness;
Palace wins at scale on both wall clock and memory. (Palace at 16 ranks
on hyperthreaded vCPUs was refused by MPI process binding; we exclude it
rather than report a degraded configuration.) The once-planned
matrix-free interior eigensolve is no longer offered as the scale
answer: its deep-shift inner solve is the measured wall of
Section~\ref{sec:roadmap}, and the honest scale story for the direct
path is the one in Table~\ref{tab:cpu}.

\subsection{GPU cell: correctness proven, scaling an honest negative}
\label{sec:perf-gpu}

geode-fem's GPU path accelerates the \emph{driven} solve (matrix-free
apply + on-device COCG), not the eigensolve of the credential, and all
GPU execution is CUDA-f32 (Section~\ref{sec:architecture}). On
2026-07-14, geode-fem code executed on a physical GPU for the first time
(EC2 \texttt{g6e.xlarge}, 1$\times$ NVIDIA L40S 46\,GB, driver
595.71.05, CUDA 13.2): the two CUDA-f32 correctness smokes --- the
matrix-free N\'ed\'elec matvec and the on-device COCG --- pass against
the f64 CPU reference within f32-appropriate tolerances ($15.0$ and
$3.3$\,s wall, durations dominated by GPU init and kernel JIT ---
correctness evidence, not performance).

The scaling study (committed as
\path{benchmarks/gpu_driven_scaling/results.toml}; a $\sigma$-lossy
lumped-port cube fixture at fixed frequency, mesh-refined from $1{,}854$
to $25{,}695$ edges, all sixteen cells on the same host) is an honest
negative: \textbf{GPU-f32 matrix-free loses to every CPU configuration
at every measured size} (Table~\ref{tab:gpu}). Computed from the
artifact's unrounded medians, the assembled-CSR COCG on CPU is
$136\times$ faster at $1{,}854$ edges (displayed: $0.032$ versus
$4.39$\,s) and still $44\times$
faster at $25{,}695$ ($1.86$ versus $81.8$\,s); even the direct LU is
$13.5\times$ faster at the top size ($6.04$ versus $81.8$\,s). The
mechanism is kernel-launch domination (${\sim}28$\,ms per iteration at
the top size versus ${\sim}0.63$\,ms for assembled-CSR on CPU) ---
$25.7$k edges is far too small to fill an L40S --- and f32 accuracy
compounds the verdict (relative error degrading to $1.2\times10^{-3}$
with refinement; nondeterministic convergence at the top size). One
directional positive survives, and it is the one the architecture
predicts: against the \emph{same matrix-free algorithm on CPU}, the
GPU's deficit falls monotonically to parity at ${\sim}26$k edges
($81.8$ versus $82.1$\,s). We state the GPU position as trajectory, not
achievement, with tracked levers (kernel fusion; on-device
f64/mixed-precision when cubecl ships it; preconditioning) --- and we
prefer publishing the bound to implying the win. This cell is a cube
fixture family, not the transmon eigenproblem, and its CPU baselines
(4-vCPU host) are not comparable to the \texttt{m6i} numbers of
Table~\ref{tab:cpu}.

\begin{table}[t]
  \centering
  \caption{Driven-solve wall-clock scaling, solve-only medians of $n=3$
  timed repetitions (warm-up excluded), same host (\texttt{g6e.xlarge}:
  4 vCPU + 1$\times$ L40S). One fixture ($\sigma$-lossy lumped-port
  cube) at fixed $\omega$, mesh-refined across the columns. f64
  iterative configurations at tolerance $10^{-8}$; the f32 GPU
  configuration at $10^{-6}$ (its recurrence floor sits above
  $10^{-8}$). The accuracy row is the full-field relative $L_2$ error of
  the GPU-f32 solution against the Direct-f64 reference. At 25{,}695
  edges the GPU-f32 configuration converged in the committed
  single-frequency runs but not in its five-point frequency sweep, and
  run-to-run convergence is nondeterministic at that size (f32 floor).
  Values verbatim from the committed
  \texttt{benchmarks/gpu\_driven\_scaling/results.toml}.}
  \label{tab:gpu}
  \small
  \setlength{\tabcolsep}{5pt}
  \begin{tabular}{@{}lcccc@{}}
    \toprule
    & \multicolumn{4}{c}{Mesh size (N\'ed\'elec edges)} \\
    \cmidrule(l){2-5}
    Configuration & 1{,}854 & 5{,}859 & 13{,}428 & 25{,}695 \\
    \midrule
    Direct faer LU (CPU f64), s & 0.024 & 0.203 & 1.540 & 6.036 \\
    Assembled-CSR COCG (CPU f64), s & 0.032 & 0.198 & 0.709 & 1.865 \\
    Matrix-free COCG (CPU f64), s & 1.652 & 8.885 & 30.234 & 82.056 \\
    Matrix-free COCG (GPU f32), s & 4.388 & 13.351 & 34.381 & 81.764 \\
    \midrule
    GPU-f32 accuracy vs.\ Direct-f64 (rel.\ $L_2$) &
      $1.26\times10^{-4}$ & $2.93\times10^{-4}$ &
      $4.60\times10^{-4}$ & $1.20\times10^{-3}$ \\
    \bottomrule
  \end{tabular}
\end{table}

\section{Reproducibility}
\label{sec:repro}

Every artifact behind every number in this paper is committed and
scripted.\footnote{Repository URL and archival DOI for the geode-fem
source to be added at submission. TODO(operator).}

\begin{itemize}
  \item \textbf{Geometry and mesh.} Generated by DeviceLayout.jl
    v1.15.0's \texttt{SingleTransmon} example \citep{devicelayout}; the
    fixture generator is committed, and the MSH~4.1 fixture is pinned by
    SHA-256 (\texttt{5b3ff4c3\ldots{}b33dd}; full hash and provenance in
    the committed \texttt{transmon\_smoke.provenance.txt}).
  \item \textbf{The MSH conversion gotcha.} MFEM's MSH~4.1 reader
    rejects the DeviceLayout multi-entity-block file, so the Palace runs
    consume an MSH~2.2 conversion (\texttt{gmsh -save -format msh2}).
    The conversion is node-preserving (all 22{,}684 nodes unchanged,
    physical groups preserved) and physics-neutral: eigenvalues
    reproduce bit-for-bit across the conversion. Documented in the
    committed \texttt{palace\_config.provenance.txt} because it is
    exactly the kind of silent pipeline step a same-mesh claim must
    disclose.
  \item \textbf{Credential artifacts.} The Palace configuration and
    provenance under \path{reference/fixtures/transmon_palace/};
    geode-fem's benchmark definitions, the agreement numbers, the
    L-doubling tripwire ($17.4901 \to 12.37$\,GHz, committed ratio
    $0.7071$), the spurious-mode measurements ($3.4528$\,GHz,
    $p = 0.9942$; L-doubling ratio $0.7081$), and the full
    gauge/projection arc of Section~\ref{sec:spurious-mode}
    (tree--cotree counts and the $1.64\%$ shift; bulk-projection
    divergence ratio $50.2$ and projected norm $1.06\times10^{-4}$;
    port-aware acceptance at $\leq 0.029\%$) in
    \path{benchmarks/transmon_eigen/results.toml}, each pinned by a
    release-tier regression test. Palace's raw outputs
    (\texttt{eig.csv}, per-port participation, run log) are committed;
    a rerun reproduced \texttt{eig.csv} bit-for-bit.
  \item \textbf{Precision of the committed agreement numbers.}
    \path{benchmarks/transmon_eigen/results.toml} stores geode-fem
    frequencies at three decimals while the Palace column carries full
    precision, so the committed per-mode \texttt{rel\_err\_pct} is
    computed against the rounded geode-fem values; at full precision the
    agreement is slightly better (approximately $0.029\%$ worst case).
    We quote the committed-artifact numbers throughout and disclose the
    rounding here. Derived ratios quoted in the text (core-second and
    speedup multipliers) are computed from the displayed table values
    unless marked as committed.
  \item \textbf{LOM chain.} The three-conductor capacitance extraction
    and Koch-exact spectrum ($C_\Sigma = 136.7$\,fF, $E_C =
    0.142$\,GHz, and the anchor-gap diagnosis of
    Section~\ref{sec:lom}) in
    \path{benchmarks/transmon_quantum/results.toml}; pinned by the
    release test \path{crates/geode-core/tests/transmon_quantum.rs}.
  \item \textbf{Adjoint layer.} The four sensitivity modules with their
    FD-and-mutation test suites:
    \path{crates/geode-core/src/adjoint.rs} (PR \#573);
    \path{crates/geode-core/src/shape.rs} (PRs \#575/\#586; this module
    also carries \texttt{capacitance\_shape\_gradient});
    \path{crates/geode-core/src/driven/adjoint.rs} (PR \#579); and
    \path{crates/geode-core/src/driven/shape.rs} (PR \#581).
  \item \textbf{Optimization demonstrations.} The parallel-plate loop
    proof (full trajectories, closed-form check, fresh-solve
    confirmation) in \path{benchmarks/transmon_diffopt/results.toml};
    the real-device FD validation, demonstration convergence, and the
    bounded anchor attempt with its mesh-safety budget in
    \path{benchmarks/transmon_diffopt/pad_results.toml}; both consumed
    by \path{crates/geode-core/tests/transmon_diffopt.rs}.
  \item \textbf{Performance artifacts.} The matched CPU cell in
    \path{benchmarks/transmon_bench_cpu/results.toml}; the 1.16M-DOF
    completion run log at
    \path{benchmarks/transmon_bench_cpu/geode_runs_1p16M_2026-07-15.log};
    the GPU scaling cell (all sixteen cells, iteration counts,
    recomputed true residuals, DNF handling, same-host provenance) in
    \path{benchmarks/gpu_driven_scaling/results.toml}; the deep-shift
    eigensolve characterization at
    \path{benchmarks/transmon_bench_cpu/sigma4p5_deepshift_characterization.md}.
  \item \textbf{Hardware disclosure.} CPU cell: EC2 \texttt{m6i.4xlarge}
    (8 physical cores / 16 vCPU, 64\,GB), \texttt{us-west-2}; the
    1.16M-DOF completion row: a 128\,GB \texttt{r6i}-class box (see the
    committed log). GPU cell: EC2 \texttt{g6e.xlarge} (4 vCPU, 30\,GB;
    1$\times$ NVIDIA L40S 46\,GB; driver 595.71.05, CUDA 13.2).
  \item \textbf{Versions.} geode-fem commit \texttt{3174015} (headline
    credential and CPU cells; the adjoint/diffopt layers are merged
    post-\texttt{3174015} and cited by PR number); Palace commit
    \texttt{fba6a5b}; DeviceLayout.jl v1.15.0; Burn/burn-cuda 0.21.
\end{itemize}

\section{Discussion}
\label{sec:discussion}

\paragraph{What the credential does and does not establish.}
In V\&V terms \citep{roache1997quantification,oberkampf2010verification}
the cross-validation is code verification and code-to-code comparison,
not validation against nature: two independent implementations agreeing
to $\leq 0.033\%$ on the identical discrete problem, with analytic
tripwires pinning the junction physics. It does not certify that either
solver matches a measured device --- that is what
\citet{ye2025electromagnetic}-style experimental benchmarks address ---
but it rules out the large error class living in formulation, assembly,
boundary handling, and eigensolver machinery. For this paper's purpose
it establishes the specific premise the gradients need: the forward
solutions being differentiated are production-accurate.

\paragraph{What the differentiable-design evidence supports.}
The proven capability is bounded and we state its boundary. Proven: the
$2\times 2$ adjoint matrix at operator level (FD-validated with mutation
tests); the capacitance$\to E_C$ chain with its stationarity collapse;
a closed Newton loop confirmed by independent fresh solves on both an
analytic fixture and the real 133k-tet mesh; and a real-mesh gradient
validated at $1.15\times10^{-4}$ against the full independent pipeline.
Not proven, and not claimed: eigenmode/EPR derivatives (the
Section~\ref{sec:roadmap} wall); anchor-reaching real-device
optimization (the $33\times$ parametrization shortfall of
Section~\ref{sec:optdemo-pad}); tensor-$\varepsilon$ shape chains;
multi-parameter optimization at scale; and any wall-clock superiority
over sweep-based workflows. The capability is also not qubit-specific
--- shape and material sensitivities of electrostatic and driven-EM
observables apply to RF and photonic device design generally; the
transmon is the demonstration vehicle with the clearest documented
unmet need.

\paragraph{Threats to validity.}
A common-mode error --- both solvers wrong the same way --- cannot be
excluded by cross-comparison alone; the analytic anchors and disjoint
implementations make it unlikely for the junction physics. The
optimization demonstrations are single-parameter by design; nothing here
demonstrates multi-parameter or constrained optimization behavior. The
parallel-plate tier's one-step convergence is an affine-map property,
disclosed as such. The real-device chain uses the scalar trace-averaged
sapphire (tensor-$\varepsilon$ delta quantified at $0.75\%$).
Performance numbers are $n=3$ measurements on one instance type per
cell and support architecture-level readings, not microsecond claims;
the large-scale completion row is a different (128\,GB) box, marked as
such in Table~\ref{tab:cpu}.

\paragraph{Limitations.}
(i) One geometry, one mesh, first order only --- higher-order and
mesh-refinement convergence is future work (a $p = 2$ Palace
configuration is committed for that purpose). (ii) The lossless model
leaves readout ports open; resistive terminations (hence
complex-symmetric pencils and mode $Q$s) are out of scope. (iii) The
credential eigenpath remains the ungauged committed benchmark; the
port-aware projection retains all six modes at $\leq 0.029\%$ but is a
composite solver, and the solenoidal port artifact is characterized and
filtered rather than eliminated. (iv) The CPU per-core win is
small-to-medium-scale only: at ${\sim}1.16$M DOFs the direct solve
completes given memory but loses to the distributed iterative reference
on both wall clock and memory --- a flop-and-fill crossover below 1M
DOFs, per the committed 128\,GB-box measurement. (v) GPU evidence covers
driven-solve correctness and an honest scaling negative in f32 on a
cube fixture family; no GPU eigensolve path exists. (vi) The
differentiable branch is electrostatic/LOM only (eigenmode/EPR is
roadmap), currently scalar-$\varepsilon$ for the shape chain, and its
real-device demonstration is bounded by the fixed-topology
parametrization budget.

\section{Conclusion}
\label{sec:conclusion}

Superconducting-qubit electromagnetic design iterates a solver that
cannot say which way to move. This paper demonstrated the missing
derivative: a discrete-adjoint layer over a tensor-native FEM solver
that turns one extra solve per step into exact, solver-derived
sensitivities --- a validated $2\times 2$ matrix of material and
geometry gradients on the electrostatic and driven $H(\mathrm{curl})$
operators --- composed into a differentiable capacitance$\to E_C$ chain
whose shape derivative collapses, by variational stationarity, to a
Hellmann--Feynman-like explicit term. The capability was exercised
honestly at two tiers: a closed Newton loop that hits a $0.2156$\,GHz
target $E_C$ and survives an independent fresh-solve check at
$1.4\times10^{-15}$ (one-step convergence disclosed as the affine-map
property it is), and the real 133k-tet device mesh, where the gradient
is FD-validated at $1.15\times10^{-4}$, a within-budget target converges
in two genuine Newton steps --- and the ${\sim}90$\,fF design anchor
sits $33\times$ beyond the fixed-topology parametrization's mesh budget,
a mesh-morphing problem we publish as the capability's current boundary.
The claim is held to a cross-validation evidence standard: the same
solver reproduces Palace to $\leq 0.033\%$ across all six eigenmodes of
the identical mesh --- a credential, not the contribution, and
incidentally the first independent validation of the reference stack.
Around the capability we kept the honest ledger: no measured performance
corner over Palace at scale (the direct eigensolve's advantage is
small-to-medium-bounded, with a flop-and-fill crossover below 1M DOFs),
a GPU path that is correctness-proven and performance-negative at
benchmark sizes, and an eigenmode/EPR branch that is roadmap --- blocked
at a characterized deep-shift eigensolve wall, with the
Jacobi--Davidson-plus-Helmholtz-projection path and the adjoint-eigenpair
formulas named. The position we defend is the one the evidence supports:
gradient-based electromagnetic design of the transmon's electrostatic
Hamiltonian parameters works today, from the field solve itself, on
production geometry --- and the boundary of that capability is published
alongside it.

\paragraph{Acknowledgments.}
geode-fem's development was AI-orchestrated (agent-driven implementation
with human direction and review); this benchmark program and paper
follow the same discipline, with every reported number traced to a
committed, human-auditable artifact.
% TODO(operator): finalize acknowledgment wording before submission.
% TODO(operator): whiteroom L1-L4 operator specification — decide
% public/citable (repo/DOI) vs acknowledged-not-cited (BRIEF coauthor note).

\bibliographystyle{plainnat}
\bibliography{refs}

\end{document}
